%!Mode:: "TeX:UTF-8"
%!TEX encoding = UTF-8 Unicode
%lan
%arara: xelatex
\documentclass{ctexart}
\newif\ifpreface
%\prefacetrue
\input{../../../global/all}
\begin{document}
\large
\setlength{\baselineskip}{1.2em}
\ifpreface
    \input{../../../global/preface}
\newgeometry{left=2cm,right=2cm,top=2cm,bottom=2cm}
\else
\newgeometry{left=2cm,right=2cm,top=2cm,bottom=2cm}
\maketitle
\fi
%from_here_to_type
\begin{problem}[2.1]
Compute the gradient $\nabla f(x)$ and Hessian $\nabla^2 f(x)$ of the Rosenbrock function
\[
    f(x) = 100(x_2-x_1^2)^2+(1-x_1)^2.
\]
Show that $x^*=(1,1)^T$ is the only local minimizer of this function, and that the Hessian matrix at that point is positive definite.
\end{problem}

\begin{problem}[2.5]
Consider the function $f:\mathbb{R}^2\to\mathbb{R}$ defined by
\[
    f(x)=\|x\|_2.
\]
Show that the sequence of iterates $\{x_k\}$ defined by
\[
    x_k=
    \left(1+\frac{1}{2^k}\right)
    \begin{pmatrix}
        \cos k\\
        \sin k
    \end{pmatrix}
\]
satisfies
\[
    f(x_{k+1})<f(x_k)
\]
for $k=0,1,2,\ldots$. Show that every point on the unit circle
\[
    \{x\mid \|x\|_2=1\}
\]
is a limit point for $\{x_k\}$.

Hint: Every value $\theta\in[0,2\pi]$ is a limit point of the subsequence
$\{\xi_k\}$ defined by
\[
    \xi_k
    =k\pmod{2\pi}
    =k-2\pi\left\lfloor\frac{k}{2\pi}\right\rfloor,
\]
where the operator $\lfloor\cdot\rfloor$ denotes rounding down to the next integer.
\end{problem}

\begin{problem}[2.6]
Prove that all isolated local minimizers are strict.

Hint: Take an isolated local minimizer $x^*$ and a neighborhood $\mathcal{N}$.
Show that for any $x\in\mathcal{N}$, $x\neq x^*$, we must have
\[
    f(x)>f(x^*).
\]
\end{problem}

\begin{problem}[2.7]
Suppose that
\[
    f(x)=x^TQx,
\]
where $Q$ is an $n\times n$ symmetric positive semidefinite matrix.
Show using the definition (1.4) that $f(x)$ is convex on the domain $\mathbb{R}^n$.

Hint: It may be convenient to prove the following equivalent inequality:
\[
    f(y+\alpha(x-y))
    -\alpha f(x)
    -(1-\alpha)f(y)
    \leq 0,
\]
for all $\alpha\in[0,1]$ and all $x,y\in\mathbb{R}^n$.
\end{problem}

\begin{problem}[2.9]
Consider the function
\[
    f(x_1,x_2)=(x_1+x_2^2)^2.
\]
At the point $x^T=(1,0)$ we consider the search direction
\[
    p^T=(-1,1).
\]
Show that $p$ is a descent direction and find all minimizers of the problem (2.10).
\end{problem}

\begin{problem}[2.13]
(For this and the following three questions, refer to the material on
``Rates of Convergence'' in Section A.2 of the Appendix.)

Show that the sequence
\[
    x_k=\frac{1}{k}
\]
is not Q-linearly convergent, though it does converge to zero.
(This is called sublinear convergence.)
\end{problem}

\begin{problem}[2.14]
Show that the sequence
\[
    x_k=1+(0.5)^{2^k}
\]
is Q-quadratically convergent to $1$.
\end{problem}

\begin{problem}[2.15]
Does the sequence
\[
    x_k=\frac{1}{k!}
\]
converge Q-superlinearly? Q-quadratically?
\end{problem}

\begin{problem}[2.16]
Consider the sequence $\{x_k\}$ defined by
\[
    x_k=
    \begin{cases}
        \left(\dfrac14\right)^{2^k}, & k \text{ even},\\[2mm]
        x_{k-1}/k, & k \text{ odd}.
    \end{cases}
\]
Is this sequence Q-superlinearly convergent? Q-quadratically convergent?
R-quadratically convergent?
\end{problem}


\begin{problem}[3.1]
Program the steepest descent and Newton algorithms using the backtracking line
search, Algorithm 3.1. Use them to minimize the Rosenbrock function (2.22).
Set the initial step length $\alpha_0=1$ and print the step length used by each
method at each iteration. First try the initial point
\[
    x_0=(1.2,1.2)^T
\]
and then the more difficult starting point
\[
    x_0=(-1.2,1)^T.
\]
\end{problem}

\begin{problem}[3.4]
Show that the one-dimensional minimizer of a strongly convex quadratic function
always satisfies the Goldstein conditions (3.11).
\end{problem}

\begin{problem}[3.5]
Prove that
\[
    \|Bx\|\geq\frac{\|x\|}{\|B^{-1}\|}
\]
for any nonsingular matrix $B$. Use this fact to establish (3.19).
\end{problem}

\begin{problem}[3.7]
Prove the result (3.28) by working through the following steps.

First, use (3.26) to show that
\[
    \|x_k-x^*\|_Q^2-\|x_{k+1}-x^*\|_Q^2
    =
    2\alpha_k\nabla f_k^TQ(x_k-x^*)
    -\alpha_k^2\nabla f_k^TQ\nabla f_k,
\]
where $\|\cdot\|_Q$ is defined by (3.27).

Second, use the fact that
\[
    \nabla f_k=Q(x_k-x^*)
\]
to obtain
\[
    \|x_k-x^*\|_Q^2-\|x_{k+1}-x^*\|_Q^2
    =
    2\frac{(\nabla f_k^T\nabla f_k)^2}
            {\nabla f_k^TQ\nabla f_k}
    -
    \frac{(\nabla f_k^T\nabla f_k)^2}
         {\nabla f_k^TQ\nabla f_k}
\]
and
\[
    \|x_k-x^*\|_Q^2
    =
    \nabla f_k^TQ^{-1}\nabla f_k.
\]
\end{problem}

\begin{problem}[3.8]
Let $Q$ be a positive definite symmetric matrix. Prove that for any vector $x$,
we have
\[
    \frac{(x^Tx)^2}
         {(x^TQx)(x^TQ^{-1}x)}
    \geq
    \frac{4\lambda_n\lambda_1}
         {(\lambda_n+\lambda_1)^2},
\]
where $\lambda_n$ and $\lambda_1$ are, respectively, the largest and smallest
eigenvalues of $Q$.

(This relation, which is known as the Kantorovich inequality, can be used to
deduce (3.29) from (3.28).)
\end{problem}

\begin{problem}[3.9]
Program the BFGS algorithm using the line search algorithm described in this
chapter that implements the strong Wolfe conditions. Have the code verify that
\[
    y_k^Ts_k
\]
is always positive. Use it to minimize the Rosenbrock function using the
starting points given in Exercise 3.1.
\end{problem}

\begin{problem}[3.12]
Consider a block diagonal matrix $B$ with $1\times1$ and $2\times2$ blocks.
Show that the eigenvalues and eigenvectors of $B$ can be obtained by computing
the spectral decomposition of each diagonal block separately.
\end{problem}
\end{document}
